% ══════════════════════════════════════════════════════════════════════════════
% Lecture 17: Union-Find and Minimum Spanning Trees
% ══════════════════════════════════════════════════════════════════════════════

% ══════════════════════════════════════════════════════════════════════════════
\section*{Disjoint-Set Data Structures (Union-Find)}
\addcontentsline{toc}{subsection}{Union-Find Data Structure}
% ══════════════════════════════════════════════════════════════════════════════

% ─────────────────────────────────────────────────────────────────────────────
\subsection{Introduction}

\begin{definition}{Disjoint-Set Data Structure}{disjointset}
A \textbf{disjoint-set data structure} (also known as \textbf{union-find}) maintains a collection 
$\mathcal{S} = \{S_1, S_2, \ldots, S_k\}$ of disjoint dynamic sets. Each set is identified by 
a \textbf{representative}, which is some member of the set. The representative doesn't matter 
as long as querying the representative twice without modifying the set returns the same answer.
\end{definition}

The union-find data structure supports three fundamental operations:

\begin{itemize}[noitemsep]
  \item \textbf{Make-Set}$(x)$: Create a new set $S_i = \{x\}$ and add $S_i$ to $\mathcal{S}$.
  \item \textbf{Union}$(x, y)$: If $x \in S_x$ and $y \in S_y$, replace $S_x$ and $S_y$ with 
        $S_x \cup S_y$ in $\mathcal{S}$. The representative of the new set is any member in 
        $S_x \cup S_y$ (often the representative of one of $S_x$ or $S_y$). This operation 
        destroys $S_x$ and $S_y$ since sets must remain disjoint.
  \item \textbf{Find}$(x)$: Return the representative of the set containing $x$.
\end{itemize}

% ─────────────────────────────────────────────────────────────────────────────
\subsection{Application: Connected Components}

A classic application of union-find is computing connected components in an undirected graph 
$G = (V, E)$. Two vertices $u$ and $v$ are in the same connected component if and only if 
there is a path between them. Connected components partition vertices into equivalence classes.

The algorithm proceeds as follows:
\begin{itemize}[noitemsep]
  \item Initialize: For each vertex $v \in V$, call Make-Set$(v)$.
  \item For each edge $(u, v) \in E$: if Find$(u) \neq$ Find$(v)$, call Union$(u, v)$.
\end{itemize}

This requires $|V|$ Make-Set operations, at most $2|E|$ Find operations, and at most $|V|-1$ 
Union operations, for a total of at most $|V| + 3|E|$ operations.

% ─────────────────────────────────────────────────────────────────────────────
\subsection{Linked List Representation}

One implementation represents each set as a singly linked list:
\begin{itemize}[noitemsep]
  \item Each set object maintains a pointer to the head (the representative) and the tail.
  \item Each list element contains the set member, a pointer to the next element, and a 
        pointer back to the set object.
\end{itemize}

\textbf{Operations:}
\begin{itemize}[noitemsep]
  \item Make-Set$(x)$: Create a singleton list in time $\Theta(1)$.
  \item Find$(x)$: Follow the pointer to the set object, then to the representative in $\Theta(1)$.
  \item Union$(x, y)$: Append $y$'s list onto the end of $x$'s list using $x$'s tail pointer. 
        Update the set pointer for every element in $y$'s list. This can take $\Theta(n)$ time 
        in the worst case.
\end{itemize}

\textbf{Weighted-Union Heuristic:} Always append the smaller list to the larger list (break ties arbitrarily).

\begin{theorem}{Linked List with Weighted-Union}{linkedlistuf}
With the weighted-union heuristic, a sequence of $m$ operations (Make-Set, Union, Find) on 
$n$ elements takes $O(m + n \log n)$ time.
\end{theorem}

\begin{proof}
Make-Set and Find still take constant time. The key observation is: how many times can each 
element's representative pointer be updated? Each time an element's pointer is updated, it must 
be in the smaller set being appended to a larger set. Thus the size of the set containing that 
element at least doubles. Since the maximum set size is $n$, each element's pointer can be 
updated at most $\log n$ times. With $n$ elements and at most $n$ unions, the total cost of 
all pointer updates is $O(n \log n)$.
\end{proof}

For the connected components problem with $|V|$ elements and $\le |V| + 3|E|$ operations, 
the total running time is $O(|V| \log |V| + |E|)$.

% ─────────────────────────────────────────────────────────────────────────────
\subsection{Disjoint-Set Forest}

A more efficient implementation represents each set as a rooted tree:
\begin{itemize}[noitemsep]
  \item Each node $x$ has a parent pointer $x.p$.
  \item The root of each tree is the representative of the set.
  \item Each node also maintains a rank $x.\text{rank}$ (an upper bound on the height).
\end{itemize}

\textbf{Operations:}

\begin{algorithm}[H]
\caption{Make-Set and Find-Set}
\begin{algorithmic}[1]
\Procedure{Make-Set}{$x$}
  \State $x.p \gets x$
  \State $x.\text{rank} \gets 0$
\EndProcedure
\end{algorithmic}
\begin{algorithmic}[1]
\Procedure{Find-Set}{$x$}
  \If{$x \neq x.p$}
    \State $x.p \gets$ \Call{Find-Set}{$x.p$} \AlgComment{Path compression}
  \EndIf
  \State \Return $x.p$
\EndProcedure
\end{algorithmic}
\end{algorithm}

\begin{algorithm}[H]
\caption{Union by Rank}
\begin{algorithmic}[1]
\Procedure{Union}{$x, y$}
\State \Call{Link}{\text{Find-Set}(x), \text{Find-Set}(y)}
\EndProcedure

\Statex % Adds a small gap between procedures

\Procedure{Link}{$x, y$}
    \If{$x.\text{rank} > y.\text{rank}$}
        \State $y.p \gets x$
    \Else
        \State $x.p \gets y$
        \If{$x.\text{rank} = y.\text{rank}$}
            \State $y.\text{rank} \gets y.\text{rank} + 1$
        \EndIf
    \EndIf
\EndProcedure
\end{algorithmic}
\end{algorithm}

\textbf{Union by rank:} Always attach the tree with smaller rank under the root of the tree 
with larger rank.

\textbf{Path compression:} During Find-Set$(x)$, make each node on the path from $x$ to the 
root point directly to the root.

\begin{complexity}{Forest with Union-by-Rank and Path Compression}{forestuf}
If we use both union by rank and path compression, a sequence of $m$ operations (Make-Set, 
Union, Find) on $n$ elements takes $O(m \cdot \alpha(n))$ time, where $\alpha(n)$ is the 
\textbf{inverse Ackermann function}.
\end{complexity}

\begin{remark}{Inverse Ackermann Function}{invack}
The function $\alpha(n)$ is an extremely slowly growing function. For all practical purposes, 
$\alpha(n) \le 4$:
\begin{center}
\begin{tabular}{cc}
\toprule
$n$ & $\alpha(n)$ \\
\midrule
$0$ to $2$ & $0$ \\
$3$ & $1$ \\
$4$ to $7$ & $2$ \\
$8$ to $2047$ & $3$ \\
$2048$ to $A(4,2)$ & $4$ \\
\bottomrule
\end{tabular}
\end{center}
where $A(4,2)$ is a number with more than $10^{19000}$ digits. Thus, for the connected 
components problem, the running time is $O((|V| + |E|) \alpha(|V|)) \approx O(|V| + |E|)$, 
which is essentially linear.
\end{remark}

% ══════════════════════════════════════════════════════════════════════════════
\section*{Minimum Spanning Trees}
\addcontentsline{toc}{subsection}{Minimum Spanning Trees}
% ══════════════════════════════════════════════════════════════════════════════

% ─────────────────────────────────────────────────────────────────────────────
\subsection{Problem Definition}

\begin{definition}{Minimum Spanning Tree (MST)}{mst}
\textbf{Input:} An undirected, connected graph $G = (V, E)$ with a weight function 
$w: E \to \mathbb{R}$ that assigns a weight $w(u,v)$ to each edge $(u,v) \in E$.

\textbf{Output:} A spanning tree $T \subseteq E$ of minimum total weight 
$w(T) = \sum_{(u,v) \in T} w(u,v)$.

A \textbf{spanning tree} is an acyclic subset of edges that connects all vertices. 
It has exactly $|V| - 1$ edges.
\end{definition}

% ─────────────────────────────────────────────────────────────────────────────
\subsection{Applications}

\textbf{Communication Networks:} A multinational company wants to lease communication lines 
between its various locations. Each pair of locations has a cost for leasing a direct line. 
The company wants to connect all locations with minimum total cost. The solution is the 
minimum spanning tree of the graph where vertices are locations and edge weights are 
leasing costs.

% ─────────────────────────────────────────────────────────────────────────────
\subsection{Generic Greedy Approach}

Both Kruskal's and Prim's algorithms are based on a generic greedy strategy. The key 
theoretical tool is the \textbf{cut property}.

\begin{definition}{Cut}{cut}
A \textbf{cut} $(S, V \setminus S)$ is a partition of the vertices into two disjoint sets 
$S$ and $V \setminus S$. An edge $(u,v)$ \textbf{crosses} the cut if one endpoint is in $S$ 
and the other is in $V \setminus S$.
\end{definition}

\begin{theorem}{Cut Property (Safe Edge Theorem)}{cutprop}
Consider a cut $(S, V \setminus S)$ and let:
\begin{itemize}[noitemsep]
  \item $T$ be a tree on $S$ which is part of some MST
  \item $e$ be a crossing edge of minimum weight
\end{itemize}
Then there exists an MST of $G$ containing both $e$ and $T$.
\end{theorem}

\begin{proof}
If $e$ is already in the MST, we are done. Otherwise, add $e$ to the MST. This creates a 
cycle with at least one other crossing edge $f$ (since the MST was already spanning). 
We have $w(f) \ge w(e)$ since $e$ is a minimum weight crossing edge. Replace $f$ by $e$ 
in the MST. The total weight doesn't increase: $w(\text{new MST}) = w(\text{old MST}) - w(f) + w(e) \le w(\text{old MST})$. 
Thus the new tree is also an MST, and it contains both $T$ and $e$.
\end{proof}

% ─────────────────────────────────────────────────────────────────────────────
\subsection{Prim's Algorithm}

\textbf{Idea:} Start with an arbitrary vertex $v$ and set $T = \{v\}$. Greedily grow the 
tree $T$: at each step, add to $T$ a minimum-weight crossing edge with respect to the cut 
$(T, V \setminus T)$.

\textbf{Correctness:} At each iteration, by the cut property, the minimum crossing edge is 
safe to add. The algorithm terminates when $T$ is a spanning tree.

\textbf{Implementation Challenge:} How do we find the minimum crossing edge at every iteration?

\begin{itemize}[noitemsep]
  \item \textbf{Naive approach:} Check all outgoing edges. This requires $O(|E|)$ comparisons 
        at each iteration, giving $O(|E| \cdot |V|)$ total time.
  \item \textbf{Efficient approach:} For every vertex $w \notin T$, maintain $\text{dist}(w)$, 
        the weight of the lightest edge from $w$ to any vertex in $T$. When a new vertex $u$ 
        is added to $T$, check whether neighbors of $u$ decrease their distance to $T$. Use a 
        min-priority queue to extract the vertex with minimum distance.
\end{itemize}

\begin{complexity}{Prim's Algorithm}{prim}
With a binary min-heap, Prim's algorithm runs in $O(|E| \log |V|)$ time:
\begin{itemize}[noitemsep]
  \item $|V|$ Extract-Min operations: $O(|V| \log |V|)$
  \item At most $|E|$ Decrease-Key operations: $O(|E| \log |V|)$
\end{itemize}
With a Fibonacci heap, the time improves to $O(|E| + |V| \log |V|)$.
\end{complexity}

% ─────────────────────────────────────────────────────────────────────────────
\subsection{Kruskal's Algorithm}

\textbf{Idea:} Start with an empty forest $T$. Sort all edges by weight in non-decreasing 
order. Greedily add the cheapest edge that does not create a cycle.

\begin{algorithm}[H]
\caption{Kruskal's MST Algorithm}
\begin{algorithmic}[1]
\Procedure{MST-Kruskal}{$G, w$}
  \State $A \gets \emptyset$
  \For{each vertex $v \in V$}
    \State \Call{Make-Set}{$v$} \AlgComment{Initialize disjoint sets}
  \EndFor
  \State Sort edges of $E$ by weight in non-decreasing order
  \For{each edge $(u, v) \in E$ in sorted order}
    \If{\Call{Find-Set}{$u$} $\neq$ \Call{Find-Set}{$v$}}
      \State $A \gets A \cup \{(u, v)\}$ \AlgComment{Add edge to MST}
      \State \Call{Union}{$u, v$}
    \EndIf
  \EndFor
  \State \Return $A$
\EndProcedure
\end{algorithmic}
\end{algorithm}

\textbf{Correctness:} The algorithm maintains the invariant that $T$ is always a sub-forest 
of some MST. When we add edge $(u, v)$, the endpoints are in different connected components. 
Consider the cut $(S, V \setminus S)$ where $S$ is the component containing $u$. Since $(u, v)$ 
is the cheapest remaining edge and it crosses this cut, by the cut property, it is safe to add.

\textbf{Implementation Challenge:} We need to check whether adding edge $(u, v)$ creates a 
cycle. This is equivalent to checking whether $u$ and $v$ are already in the same connected 
component. This is exactly what the union-find data structure provides.

\begin{complexity}{Kruskal's Algorithm}{kruskal}
Analysis of Kruskal's algorithm:
\begin{itemize}[noitemsep]
  \item Initialize $A$: $O(1)$
  \item First for loop: $|V|$ Make-Set operations
  \item Sort edges: $O(|E| \log |E|)$
  \item Second for loop: $O(|E|)$ Find-Set and Union operations
  \item Total union-find operations: $O((|V| + |E|) \alpha(|V|))$
\end{itemize}
Total time: $O((|V| + |E|) \alpha(|V|)) + O(|E| \log |E|) = O(|E| \log |E|) = O(|E| \log |V|)$ 
(since $|E| \le |V|^2$, we have $\log |E| = O(\log |V|)$).

If edges are already sorted, the time becomes $O(|E| \alpha(|V|))$, which is almost linear.
\end{complexity}

% ─────────────────────────────────────────────────────────────────────────────
\subsection{Lecture Summary}

\begin{tcolorbox}[colback=lightblue, colframe=keyblue, fonttitle=\bfseries,
                  title={Key Points}]
\begin{itemize}[noitemsep]
  \item \textbf{Union-Find Data Structure:} maintains disjoint sets with operations Make-Set, 
        Union, and Find. Essential for tracking connected components.
  \item \textbf{Linked list implementation:} $O(m + n \log n)$ time for $m$ operations on 
        $n$ elements with weighted-union heuristic.
  \item \textbf{Forest implementation:} $O(m \cdot \alpha(n))$ time with union-by-rank and 
        path compression, where $\alpha(n)$ is the inverse Ackermann function (practically constant).
  \item \textbf{MST Problem:} Find a spanning tree of minimum total weight in a weighted 
        undirected graph.
  \item \textbf{Cut Property:} The minimum weight crossing edge is safe to add to any partial MST.
  \item \textbf{Prim's Algorithm:} Grow the tree from a single vertex by repeatedly adding 
        the minimum crossing edge. Time: $O(|E| \log |V|)$ with binary heap.
  \item \textbf{Kruskal's Algorithm:} Sort edges and add the cheapest edge that doesn't create 
        a cycle using union-find. Time: $O(|E| \log |V|)$ (sorting dominates).
  \item Both algorithms are greedy and produce correct MSTs by the cut property.
\end{itemize}
\end{tcolorbox}